%! Mode:: "TeX:UTF-8"
%!TEX encoding = UTF-8 Unicode
%!TEX TS-program = xelatex
\documentclass{ctexart}
\usepackage{bbm}
\usepackage{tikz}
\usepackage{amsmath,amssymb,amsthm,color,mathrsfs}
\usepackage{hyperref}
\usepackage{cleveref}
\usepackage{enumitem,anysize}%
\usepackage{expl3}
\marginsize{1in}{1in}{1in}{1in}%
\theoremstyle{remark}
\newtheorem*{claim*}{Claim}
\newtheorem{lemma}{Lemma}

\DeclareMathOperator{\ot}{ordertype}
\DeclareMathOperator{\dom}{dom}
\DeclareMathOperator{\ran}{ran}
\DeclareMathOperator{\field}{field}
\DeclareMathOperator{\Ord}{Ord}

\newcommand{\cc}{\mathfrak{c}}

\def\email#1{{\texttt{#1}}}
\newcommand{\isep}[1][0pt]{\addtolength{
\itemsep}{#1}}
\newcounter{problem}
\renewcommand{\theproblem}{\Roman{problem}}
\newenvironment{problem}{\refstepcounter{problem}\noindent\color{blue}
  \begin{tikzpicture}[baseline]%
    \node at (-0.02em,0.3em) {$\mathbb{P}$};%
    \node[scale=0.7] at (0.2em,-0.0em) {R};%
    \node[scale=0.7] at (0.6em,0.4em) {O};%
    \node[scale=0.8] at (1.05em,0.25em) {B};%
    \node at (1.55em,0.3em) {L};%
    \node[scale=0.7] at (1.75em,0.45em) {E};%
    \node at (2.35em,0.3em) {M};%
\end{tikzpicture}\theproblem}{}
\crefname{problem}{Problem}{Problem}
%\renewcommand\theprob{{\Roman{problem}}}
\newcommand\mysolution{
  \begin{tikzpicture}[baseline]%
    \node at (-0.04em,0.3em) {$\mathbb{S}$};%
    \node[scale=0.7] at (0.35em,0.4em) {O};%
    \node at (0.7em,0.3em) {\textit{L}};%
    \node[scale=0.7] at (0.95em,0.4em) {U};%
    \node[scale=1.1] at (1.19em,0.32em){T};%
    \node[scale=0.85] at (1.4em,0.24em){I};%
    \node at (1.9em,0.32em){$\mathcal{O}$};%
    \node[scale=0.75] at (2.3em,0.21em){\texttt{N}};%
\end{tikzpicture}}
\newenvironment{solution}{
\begin{proof}[\mysolution]}{
\end{proof}}

%%%%%%%%
\newcommand{\calL}{\mathcal{L}}

\newcommand\<{\langle}
\renewcommand\>{\rangle}
\newcommand\eneg{\mathcal{E}_{\neg}}
\newcommand\eto{\mathcal{E}_{\to}}
\newcommand\N{\mathbb{N}}
\newcommand\subini{\subsetneqq_{init}}
\def\to{\rightarrow}
\newcommand{\calA}{\mathcal{A}}
\newcommand{\xor}{\vee}
\newcommand{\bor}{\bigvee}
\newcommand{\band}{\bigwedge}
\newcommand{\xand}{\wedge}
\newcommand{\minus}{\mathbin{\backslash}}
\newcommand{\calF}{\mathcal{F}}
\newcommand\mi[1]{\mathscr{P}(#1)}
\newcommand{\calB}{\mathcal{B}}
\newcommand{\heq}{\mathop{\hat{=}}}
\newcommand{\hneq}{\mathop{\hat{\neq}}}
\newcommand{\calR}{\mathcal{R}}
\newcommand{\hle}{\mathop{\hat{<}}}
\newcommand{\R}{\mathbb{R}}
\newcommand{\calM}{\mathcal{M}}
\newcommand{\calN}{\mathcal{N}}
\newcommand{\frA}{\mathfrak{A}}
\newcommand{\card}{\mathrm{card}}
\newcommand{\frM}{\mathfrak{M}}
\newcommand{\Th}{\mathrm{Th}}
\newcommand{\ecd}{\mathrm{EC}_\Delta}
\newcommand{\Q}{\mathbb{Q}}
\newcommand{\Z}{\mathbb{Z}}
\newcommand{\oto}{\leftrightarrow}
\newcommand{\hin}{\hat{\in}}
\newcommand{\fun}[2]{{}^{#1}#2}
\newcommand{\A}{\mathbbm{A}}

\begin{document}
\title{$\mathbb{ALGEBRAIC\ GEOMETRY}$}
\author{白永乐\\ %% 姓名
  SID: 202011150087\\ %% 学号
\email{202011150087@mail.bnu.edu.cn}} %% 电邮
\maketitle

\begin{problem}
  $X$ is a topology space. Prove that $X$ is Noetherian $\iff$ every subspace of $X$ is compact.
\end{problem}
\begin{solution}
  $\Rightarrow$: Obviously subspace of Noetherian space is Noetherian space, so we only need to prove Noetherian space is compact.
  If $X$ is not compact, then $\exists v_n\subset X$ such that $v_n$ is closed and $\bigcap_{k=1}^n v_k\neq \emptyset$ and $\bigcap_{k=1}^{\infty}v_k=\emptyset$. Then $u_n:=\bigcap_{k=1}^n v_n$ is infinitly descending chain of closed sets, and since $\bigcap_{k=1}^{\infty} u_k=\bigcap_{k=1}^{\infty} v_k=\emptyset $ so it's not finally stable.

  $\Leftarrow$:
  If $X$ is not Noetherian, assume $\{u_n\}\subset\mi{x}$ is a chain of closed sets such that $u_{n+1}\subsetneqq u_n$. Let $u:=\bigcap_{k=1}^{\infty} u_k$, consider $X\minus u\subset X$. $v_n:=u_n\minus u=u_n\cap (X\minus u)$ is closed set in $X\minus u$, but $\bigcap_{k=1}^{\infty} v_n=(\bigcap_{k=1}^{\infty} u_n)\minus u=u\minus u=\emptyset$. So $X\minus u$ is not compact, contradiction!
\end{solution}

\begin{problem}
  Given $X$ Noetherian, $A\subset X$, $A=\bigcup_{k=1}^n u_k=\bigcup_{k=1}^m v_k$ and $u_k,v_k$ is irreducible non-empty closed set, $u_i\not\subset u_j,v_i\not\subset v_j$. Prove that $m=n$, and $\exists \sigma\in S_n,\forall k\in\{1,2,\cdots n\},u_k=v_{\sigma(k)}$.
\end{problem}

\begin{lemma}
  \label{lem:ir}
  If $u$ is irreducible, $u\subset \bigcap_{k=1}^n v_k$, where $v_k$ is closed, then if $u\cap v_i\neq \emptyset$, then $u\subset v_i$. And $\exists k,u\in v_k$.
\end{lemma}
\begin{proof}
  Assume $u\cap v_1\neq \emptyset$. If $u\not\subset v_1$, then $u=(u\cap v_1)\cup(u\cap \bigcup_{k=2}^{n} v_k)$, contradiction with $u$ is irreducible. So $u\subset v_1$.

  If $u=\emptyset$ it's obvious $u\subset v_1$. If not, then $u\cap v_k,k=1,2\cdots n$ can't be all empty. So $u\subset v_k$ for some $k$.
\end{proof}

\begin{solution}
  From \Cref{lem:ir} we know $\forall i,\exists j,u_i=v_j$. If $u_i=v_j=v_k$ then $v_j\subset v_k$ thus $j=k$. So $\forall i,\exists! j,u_i=v_j$. Let $\phi:\{1,2,\cdots n\}\to\{1,2,\cdots m\},i\mapsto j$. Then $\phi$ is a map. If $\phi(i)=\phi(j)$ then $u_i=u_j$, thus $i=j$. So $\phi$ is injection. Consider $v_j$, since $v_j\subset \bigcup_{k=1}^n u_k$, so from \Cref{lem:ir} we know $\exists i,u_i=v_j$. So $\phi(i)=j$. Thus $\phi$ is bijection. So $m=n,\phi\in S_n$. Let $\sigma=\phi\in S_n$ satisfy the condition.
\end{solution}

\begin{problem}\label{pro:p}
  $K$ is a field, prove that $(xy-1)$ is prime ideal of $K[x,y]$.
\end{problem}

\begin{solution}
  Only need to prove $K[x,y]/(xy-1)$ is integral domain. Consider homeomorphism $\phi:K[x,y]\to K[x,x^{-1}],f(x,y)\mapsto f(x,x^{-1})$. Obviously $\ker(\phi)\ni xy-1$, now we prove $\ker(\phi)=(xy-1)$. $\forall f(x,y)\in \ker(\phi)$, we have $f(x,x^{-1})=0$. Obviously we can get $f(x,y)=(xy-1)g(x,y)+l(x)+m(y),g(x,y)\in K[x,y],l(x)\in K[x],m(y)\in K[y]$. So $f(x,x^{-1})=l(x)+m(x^{-1})=0$. Since $x^n,n\in \Z$ is linear  independent, so $l(x),m(x^{-1})\in K,l+m=0$. So $l(x)+m(y)=l+m=0$. So $f(x,y)\in (xy-1)$.

  So $\ker(\phi)=(xy-1)$, and thus $K(x,y)/(xy-1)\cong K[x,x^{-1}]$. So $K[x,y]/(xy-1)$ is integral domain and $(xy-1)$ is prime.
\end{solution}

\begin{problem}
  Let $k$ be a field, $I=(xy-1)$, $C=V(I)\in\mathbb{A}^2_k$.
  Show that $I(C)=(xy-1)$.
\end{problem}

\begin{solution}
  Obviously $(xy-1)\subset I(C)$, so we only need to prove $I(C)\subset (xy-1)$.
  Consider $f(x,y)\in I(C)$, we get $f(t,t^{-1})=0,\forall x\in \R\minus \{0\}$. For $n$ large enough we know $x^nf(x,x^{-1})\in \R[x]$, and it has infinite roots, so $x^nf(x,x^{-1})=0$. Thus $f(x,x^{-1})=0$. From \Cref{pro:p} we know $f\in \ker(\phi)=(xy-1)$. So $(xy-1)=I(C)$.
\end{solution}

\begin{problem}
  If $X$ is Noetherian space, $Y\subset X$, then $\dim Y\leq \dim X$.
\end{problem}

\begin{lemma}\label{lem:sub}
  If $X$ is Noetherian, $Y\subset X$, and $u\subset Y$ is irreducible closed set in $Y$, then exists least closed set $v\subset X$ such that $u\subset v$, and $v$ is irreducible in $X$.
\end{lemma}
\begin{proof}
  Let $\mathcal{W}:=\{w\subset X:u\subset w,w \text{ is irreducible closed set in }X\}$, and $v:=\bigcap_{w\in \mathcal{W}}w$, we will prove $v$ can satisfy the given condition. Obviously $v$ is closed. Assuming $v=\bigcup_{k=1}^n v_k$, where $v_k$'s are irreducible, then $u=\bigcup_{k=1}^n (v_k\cap Y)$. Since $u$ is irreducible in $Y$, we get $u\subset v_k\cap Y$ for some $k$. So $v_k\in\mathcal{W}$, and then $v\subset v_k$. So $v=v_k$ and thus irreducible. Forall closed set $t\subset X$ such that $u\subset t$, there exists a irreducible closed set $s\subset X$ such that $u\subset s\subset t$. So $s\in\mathcal{W}$ and thus $v\subset s\subset t$. So $v$ is the least.
\end{proof}

\begin{solution}
  Consider irreducible ascending chain $Y_1\subsetneqq Y_2\subsetneqq\cdots Y_n$ in $Y$. From \Cref{lem:sub} we know exists least irreducible closed $X_k\in X$ such that $Y_k\subset X_k$. Form the minimality of $X_k$ we easily get $X_k\subset X_{k+1}$. Since $Y_k$ is closed in $Y$ we get $\exists V_k\subset X$ and $V_k$ is closed and $Y_k=Y\cap V_k$. From the minimality of $X_k$ we get $X_k\subset V_k$ and thus $X_k\cap Y=Y_k$. So $X_k\cap Y\subsetneqq X_{k+1}\cap Y$, thus $X_k\subsetneqq X_{k+1}$. So $\dim X\geq \dim Y$.
\end{solution}

\end{document}


